\section{Detector}
\label{sec:detector}

%-------------------------------------------------------------------------


 \subsection{Detector Selection and Implementation}

To optimize localization speed and accuracy, we initially evaluated Google SpeciesNet \cite{speciesnet2025}, an ensemble architecture highly regarded for its robust handling of challenging environments. 
SpeciesNet utilizes the MegaDetector framework \cite{megadetector2026} and boasts an impressive 98.6\% empty image exclusion accuracy, effectively filtering out blank "wind-blows" images. 
However, while the system is highly optimized for batch processing of static, low-quality, poorly lit, or camouflaged trail-camera photographs, empirical testing revealed that its heavy computational 
overhead consistently exceeded our expectations and the strict 5-second overall inference constraint dictated by the project specifications. Consequently, as outlined in our preliminary report, 
we proceeded with YOLOv6, a lightweight and highly efficient single-stage object detector \cite{li2022yolov6}.

The selected YOLOv6 model is a standardized architecture pretrained on the comprehensive Common Objects in Context (COCO) dataset. 
This pretraining allowed us to directly leverage established COCO class identifiers (numeric IDs assigned to different object categories worldwide) to explicitly filter for our target species: 
cats (class ID 15) and dogs (class ID 16). If the detector identifies neither class, the pipeline immediately forwards the rejection class ($-1$), efficiently bypassing the downstream prediction network. 
Preliminary inference tests on our dataset demonstrated exceptional localization capabilities, achieving a 100\,\% accuracy rate in correctly identifying the primary subject, which validated this approach 
for our pipeline.

A significant concern during initial testing was the potential for grouped or closely situated animals to be merged into a single bounding box. This would severely degrade the performance of the subsequent 
fine-grained classifier, which we trained from scratch, as it would struggle to determine the correct species. However, this concern was scattered during testing with multiple complex pictures, including
 zoo animals. We observed that YOLOv6 effectively avoids this issue by creating a dedicated bounding box for each individual animal. To strictly adhere to the project's evaluation rules, our algorithm compares
  the spatial dimensions of all detected bounding boxes and forwards only the largest one.

Furthermore, to ease the representational burden on our prediction model and to eliminate the traits of other animals that may overlap with the forwarded bounding box of the searched animal,
 we implemented a zero-out masking strategy. In scenarios where smaller animals overlap with the primary subject, the overlapping regions are explicitly blacked out. We see a distinct advantage in 
 time and accuracy by eliminating the traits of the smaller animals. However, we are aware of the inherent trade-off: this masking strategy could temporarily make it harder for our model to detect 
 the primary animal, as a large chunk of its visual information might go missing.